\kbvchapter{ಪಾಠ ೬}{ಕನ್ನಡ ಸಂಧಿಗಳು}

\indent “ನಾವೆಲ್ಲರೂ ಜೊತೆಗೂಡಿ ಊರನ್ನು ಸೇರೋಣ.” - ಈ ವಾಕ್ಯದ ಅಭಿಪ್ರಾಯವೇನು ಎಂಬುದನ್ನು ನೀವು ಗ್ರಹಿಸಬಲ್ಲಿರಿ. ಆದರೆ, ಇಲ್ಲಿ ಕೆಲವು ಪದಗಳನ್ನು ಬಿಡಿಬಿಡಿಯಾಗಿ ಹೇಳಿರುವುದು ಅಷ್ಟೊಂದು ತೃಪ್ತಿಕರವಾಗಿಲ್ಲ. ವಿಚಿತ್ರ ತೆರನಾಗಿರುವ ಈ ಕ್ರಮ ರೂಢಿಗೆ ವಿರುದ್ಧವೆಂತಲೇ ಹೇಳಬಹುದು. ಈ ವಾಕ್ಯದಲ್ಲಿ ಕೆಲವೊಂದು ಕಡೆ ಅಕ್ಷರಗಳನ್ನು ಒಟ್ಟುಗೂಡಿಸಿ ಹೇಳಿ:\par

\indent - ದಲ್ಲಿ ಕೇಳಲು ಹಿತವೆನಿಸುವುದು,\par

\indent - ನಾವು + ಎಲ್ಲರೂ = ನಾವೆಲ್ಲರೂ, ಜೊತೆ + ಕೂಡಿ = ಜೊತೆಗೂಡಿ,\par

\indent - ಊರು + ಅನ್ನು = ಊರನ್ನು ಎಂದಾಗಲಿ ಸುಧಾರಿತ ವಾಕ್ಯವು “ನಾವೆಲ್ಲರೊ ಜೊತೆಗೂಡಿ ಊರನ್ನು ಸೇರೋಣ” ಎಂದಾಗುವುದು.\par

\indent ನಾವೆಲ್ಲರೊ ಎಂಬ ಪದದಲ್ಲಿ ಪೂರ್ವಾರ್ಧ ಭಾಗವು “ನಾವು” -ಮತ್ತು ಉತ್ತರಾರ್ಧ ಭಾಗವು “ಎಲ್ಲರೊ” ಆಗಿರುವುದು. ಇಲ್ಲಿ ಪೂರ್ವಾರ್ಧದ ಕಡೆಯ ಅಕ್ಷರ “ವು” ಮತ್ತು ಉತ್ತರಾರ್ಧದ ಮೊದಲ ಅಕ್ಷರ “ಎ” ಇವೆರಡೂ ಸೇರಿ “ವೆ” ಆಗಿರುವುದನ್ನು ಕಾಣಬಹುದು. ಹೀಗೆ, ಉಚ್ಚರಿಸಲು ಕಾಲ ವಿಳಂಬವಾಗದ ರೀತಿಯಲ್ಲಿ ಎರಡು ಅಕ್ಷರಗಳು ನಡಬಿಡಿ ಒಂದಕ್ಕೊಂದು ಸೇರುವುದಕ್ಕೆ “ಸಂಧಿ” ಎಂದು ಹೆಸರು.\par

\begin{kbvrule}
\noindent\kbvbullet\hspace{0.5em}ಸಂಧಿಯಾಗುವಾಗ ಒಂದು ಅಕ್ಷರ ಬಿಟ್ಟು ಹೋಗಬಹುದು. ಒಂದು ಅಕ್ಷರ ಹೊಸದಾಗಿ ಬಂದು ಸೇರಬಹುದು; ಒಂದು ಅಕ್ಷರಕ್ಕೆ ಬದಲಾಗಿ ಮತ್ತೊಂದು ಅಕ್ಷರ ಬರಬಹುದು. ಈ ರೀತಿಯಾಗಿ ಉಂಟಾಗುವ ಬದಲಾವಣೆಗಳನ್ನು ಸಂಧಿಕಾರ್ಯವೆನ್ನುತ್ತೇವೆ.
\end{kbvrule}

\vspace{6pt plus 5pt minus 1pt}\begin{center}{\kbvsecfont\fontsize{12.5pt}{16pt}\selectfont\color{kbvaccentdark} ಕನ್ನಡ ಸಂಧಿಗಳು}\end{center}\vspace{3pt plus 3pt minus 1pt}

\indent ಕನ್ನಡ ಸಂಧಿಗಳಲ್ಲಿ ಮೂರು ಪ್ರಕಾರಗಳಿವೆ; ಅವು ಯಾವುವೆಂದರೆ\par

\vspace{2pt plus 2pt}\noindent\hspace*{1.5em}(೧) ಲೋಪ ಸಂಧಿ,\hfill (೨) ಆಗಮ ಸಂಧಿ,\hfill (೩) ಆದೇಶ ಸಂಧಿ\hspace*{1.5em}\par\vspace{2pt plus 2pt}

\kbvsechead{(೧) ಲೋಪ ಸಂಧಿ:}

\indent “ನಾವೆಲ್ಲರೊ ಜೊತೆಗೂಡಿ ಮನೆಗೆ ಹೋಗೋಣ” ಎಂಬ ವಾಕ್ಯದಲ್ಲಿ ನಾವೆಲ್ಲರೊ ಎಂಬ ಸಂಧಿ ಪದವನ್ನು ಬಿಡಿಸಿ ನೋಡೋಣ.\par
